\markboth{序章}{序章}
\section*{序章：为什么 AI 能接住我们说的话？}
\addcontentsline{toc}{section}{序章：为什么 AI 能接住我们说的话？}

\begin{idea}
\noindent\emph{从一个真正值得问的问题开始}\par

你问一个语言模型：“$2+3$ 等于多少？”它回答“$5$”。
你让它解释诗句，它写出一段说明；让它修改代码，它也能提供程序。
我们真正好奇的，不是它为何能猜下一个字，
而是一套由数字运算组成的系统，为什么能够产生这样的回答。

本书就从这个问题出发。我们将先研究可计算的规则怎样学习语言模式，
再看看这些规则怎样组合成 Transformer，以及为什么一个实用 AI 助手
还需要额外训练、外部资料和工具。
\end{idea}

\subsection*{先玩一个不需要数学的预测游戏}

看到下面的句子，你会自然想到一些候选词：
\begin{center}
\begin{tcolorbox}[colback=paper,colframe=linegray,arc=2mm,boxrule=.6pt,
left=5mm,right=5mm,top=3mm,bottom=3mm]
“天突然下起大雨，他出门前拿了一把\underline{\hspace{12mm}}。”
\end{tcolorbox}
\end{center}
例如你想到“伞”。不是因为你查过固定答案表，
而是因为“下雨”“出门”“一把”共同限制了合理的后续内容。
把“下雨”换成“烈日当空”，候选内容也会随语境改变。
这说明：预测后续文字，往往必须利用上下文中的关系。

再看一个对话：
\[
\text{用户：“$2+3$ 等于多少？”}\qquad
\text{助手：“$5$。”}
\]
如果把用户提问与“助手：”放在同一段输入中，
\textbf{生成一段回答}也可以表示为在已有上下文后继续生成文字。
这为语言模型提供了一种统一的计算形式，
但能续写文字并不自动意味着会正确推理或可靠核实事实。

\subsection*{为什么偏偏练习“预测下一个 Token”？}

我们希望训练机器学习语言规律，可是为海量句子逐一撰写人工答案代价很高。
一个巧妙的办法是让现成的文字充当自己的练习题和答案：
给出前半段，让模型尝试预测后面是什么；
原始文本恰好提供了可以比较的目标。

为了让机器计算，我们先把文字切成若干\textbf{Token（文本处理单元）}，
它们包括字、词的一部分、标点等，不必与自然语言中的“词”一一对应。
本书重点学习\textbf{自回归语言模型}：
它根据已出现的 Token，给下一个 Token 的候选打出概率。
训练时，模型依赖真实文本得到误差信号并更新参数；
实际生成时，则依赖已经给出的文字和自己刚生成的 Token，
一步步继续写下去。

语言建模也有其他训练任务和架构；
真正的助手还会通过示范回答、偏好训练、检索或工具来适应用户任务。
我们选择下一 Token 预测作为贯穿教材的\textbf{小而具体的问题}，
是因为它足以把文字表示、神经网络学习、上下文读取和生成过程连在一起。

\subsection*{同一个例子，会从文字走到数学}

接下来我们会反复用到这条短句：
\[
\text{“春天”}\ \longrightarrow\ \text{“花”}\
\longrightarrow\ \text{“开了”}\ \longrightarrow\ \text{“。”}
\]
我们约定一个只有五个候选的简化词表，将这句话切为四个 Token，
编号为 $[2,4,5,3]$，另留编号 $1$ 表示“其他”。
全书的手算数据沿用这份约定：数字由我们设定，用来观察运算过程，模型中的参数则需要通过训练学习。

\chapterchecktasks{00-guide}{理解与复算}{
\item \textbf{练习题：}把“春天 花 开了 。”当作训练文本，写出每一步给模型看的前文和要预测的下一个 Token。一共有几道题？\chapterinlinehint{从第一个 Token 开始，每次多看一个 Token。}
\item \textbf{编号：}按序章的简化词表，把“花 开了 。”写成 ID 序列；编号 $1$ 留给什么？\chapterinlinehint{对照“春天、花、开了、。”的编号 $[2,4,5,3]$。}
\item \textbf{找错：}有人说：“只要模型逐 Token 生成得通顺，它提供的日期、来源和报价就必然真实。”解释逻辑漏洞。\chapterinlinehint{区分续写概率与外部证据的来源。}
}


\par\smallskip
真正的出发点是“计算机怎样学会生成有用的语言回答”。
下一 Token 预测是贯穿全书的具体训练任务。
现在要研究的第一步很朴素：
如果规则不能全靠人手写，机器能不能从例子中学出一个可调整的计算函数？
